% IEEE conference-format report. Build: latexmk -pdf main.tex
% Every number in this document must come from results/*.csv (cite the run id in a footnote).
\documentclass[conference]{IEEEtran}
\usepackage{graphicx,booktabs,amsmath,siunitx,hyperref,xcolor}
\newcommand{\tbd}[1]{\textcolor{red}{[TBD: #1]}}
\graphicspath{{figures/}}

\begin{document}
\title{A Cloud-Native 5G Testbed with an SDN Transport for Anomaly Detection and Handover Prediction}
\author{\IEEEauthorblockN{Rihem Fathallah}\IEEEauthorblockA{ENET'Com, University of Sfax, Tunisia\\rihem.fathallahengineer@gmail.com}}
\maketitle

\begin{abstract}
We build an emulated 5G standalone network (Open5GS core, srsRAN gNB and UE over a ZeroMQ virtual
radio) in a single-node Kubernetes cluster, whose N3 user plane crosses a software-defined transport
ring. Faults injected into the transport and into the network functions yield a labelled dataset of
\tbd{N} hours with \tbd{N} faults. We compare a z-score baseline, Isolation Forest and an LSTM
autoencoder at event level (recall, detection delay, false alarms per hour) and apply them to public
BGP data. Separately, an LSTM handover predictor is compared with a 3GPP A3 rule tuned per speed on
TR~38.901 signal traces. \tbd{headline results}.
\end{abstract}

\section{Introduction}
Research questions:
\begin{itemize}
\item \textbf{RQ1} On a Kubernetes-hosted Open5GS core with an SDN transport, what event recall,
false alarms per hour and median detection delay do Isolation Forest and an LSTM autoencoder reach
against a z-score baseline, per fault type and severity?
\item \textbf{RQ1b} Do the same models and features detect known anomalous events in public BGP data?
\item \textbf{RQ2} Does an LSTM predictor reduce ping-pong and radio-link-failure rates against an A3
rule tuned per speed, without raising the handover rate?
\end{itemize}
Contributions: \tbd{after results}.

\section{Testbed}
\subsection{Architecture}
\tbd{architecture figure}. Open5GS v2.7.7 (10 network functions of the 5G service-based architecture~\cite{ts23501}), srsRAN Project gNB and srsUE over
ZeroMQ, single-node k3s with the static CPU manager: the gNB, UE and UPF containers get exclusive
cores so that CPU faults cannot spill onto the real-time radio emulation or the measurement path.
\subsection{SDN transport}
N3 crosses a four-switch Open vSwitch ring (Mininet) controlled by Ryu. The controller computes a
breadth-first forwarding tree rooted at the RAN-site switch and recomputes it on every port-status
change; N3 normally follows s1--s2--s3 and is rerouted over s4 when a link on that path fails
(measured reroute: \tbd{host value}~s). A spanning-tree design was rejected: Ryu's stplib did not
reconverge after a failure on the active path.
\subsection{Monitoring}
Prometheus scrapes, every \SI{5}{s}: Open5GS NF metrics, kubelet/cAdvisor, node-exporter, the
controller's port and tree metrics, and a UE-side probe (RTT, loss, tunnel throughput, registration).

\section{Fault model and dataset}
Table~\ref{tab:faults}: fault types, injection point (one per type), mild/strong parameters.
Ground truth uses the times at which an impairment was \emph{confirmed} active and removed; the effect
window adds a \SI{60}{s} recovery grace. Windows (\SI{60}{s}, stride \SI{5}{s}) are anomalous if
$\geq 50\%$ of their samples fall in an effect window, normal if none, ambiguous otherwise (excluded
from window metrics). Splits are by time per run (60/20/20) with a \SI{60}{s} purge at boundaries.
\tbd{runs, hours, faults per type and severity, failed injections}.
\begin{table}[t]\centering\caption{Fault catalogue}\label{tab:faults}
\begin{tabular}{llll}\toprule Type & Injector & Mild & Strong\\\midrule
delay & Mininet link & \SI{10}{ms} & \SI{50}{ms}\\
loss & Mininet link & 1\,\% & 5\,\%\\
link down & Mininet link & -- & --\\
CPU & Chaos Mesh (UPF) & 1 worker, 50\,\% & 2 workers, 100\,\%\\
memory & Chaos Mesh (UPF) & 256\,MB & 768\,MB\\
pod kill & Chaos Mesh (AMF/SMF) & -- & --\\
flood & UE iperf3 & 4$\times$, 30\,Mb/s & 8$\times$, 80\,Mb/s\\\bottomrule
\end{tabular}\end{table}

\section{Anomaly detection (RQ1, RQ1b)}
Models: z-score ($|z|>3$), Isolation Forest~\cite{liu2008iforest} (threshold: best validation F1), LSTM autoencoder
(threshold: 99th percentile of validation scores on normal windows)~\cite{malhotra2016lstm}, Random Forest as a supervised
upper reference. Metrics at event level (recall, delay from $t_\text{applied}$ to the end of the first
alarm window), false alarms per hour of normal operation (alarms merged within \SI{30}{s}), and
window-level F1. No point adjustment~\cite{kim2022pa}. BGP features follow~\cite{alrousan2012bgp}. 5 seeds, mean $\pm$ std.
\tbd{Fig. detection\_summary, detection\_by\_type, detection\_delay, attribution; ablation without
link-state columns; RAN transfer; BGP table}.

\section{Handover prediction (RQ2)}
\subsection{Simulation model}
Seven omnidirectional cells (ISD \SI{500}{m}), TR~38.901~\cite{tr38901} UMa path loss (verified identical to ns-3's
implementation), spatially correlated shadowing, a smooth LOS field with soft transitions, UEs at
3/30/60/120\,km/h, \SI{40}{ms} RSRP samples with \SI{1}{dB} measurement error, TS~38.331~\cite{ts38331} L3 filter
($k=4$, coefficient rescaled to keep the \SI{200}{ms} time constant). Full-load SINR; RLF via
$Q_\text{out}=-8$\,dB, $Q_\text{in}=-6$\,dB, T310 = \SI{1}{s}; ping-pong = return within \SI{1}{s}.
\subsection{Policies and fairness}
A3 is tuned per speed over hysteresis 0--10\,dB and TTT \{40,160,320,640\}\,ms; the LSTM's threshold
and minimum time-of-stay are tuned per speed with the same objective
$J = \text{RLF/min} + 0.5\,\text{PP/min} + 0.05\,\text{HO/min}$ on the same validation trajectories.
Label: a neighbour exceeds the serving cell by $\geq 3$\,dB for $\geq 200$\,ms starting within
\SI{500}{ms} (noise-free large-scale RSRP). Training states come from A3 rollouts plus one DAgger round.
\tbd{Fig. handover\_rates, a3\_grid; paired bootstrap CIs (LSTM $-$ A3) per speed}.

\section{Limitations}
Emulated radio with a single UE (srsRAN) or protocol-level UEs (UERANSIM); synthetic faults; a
single host; simplified mobility (omnidirectional cells, no fast fading); BGP intervals taken from
the literature. Results describe the core and transport behaviour of an emulated network.


\bibliographystyle{IEEEtran}
\bibliography{refs}
\end{document}
